\documentclass[a4paper,12pt]{article}

\usepackage[T1]{fontenc}
\usepackage[utf8]{inputenc}
\usepackage[polish]{babel}
\usepackage[a4paper,margin=2.5cm]{geometry}

\usepackage{graphicx}
\usepackage{float}
\usepackage{enumitem}

\setlength{\parindent}{0pt}
\setlength{\parskip}{0.5em}

\begin{document}

\begin{center}
    {\LARGE \textbf{Sprawozdanie Projekt 2 - System planowania i obsługi powiadomień}}\\[1.5em]
\end{center}

\textbf{Numer albumu: \underline{\hspace{4cm}}}

\textbf{Imię: \underline{\hspace{5cm}}}

\textbf{Nazwisko: \underline{\hspace{4.7cm}}}

\textbf{Grupa: \underline{\hspace{5.4cm}}}

\vspace{1em}

\section*{1. Wstęp}
\subsection*{a) Opis struktury projektu - Moduł Notifications}

Moduł \texttt{Notifications} jest podzielony na warstwy:
\begin{itemize}[leftmargin=1.5em]
    \item \texttt{web/} -- endpointy API i obsługa żądań HTTP,
    \item \texttt{service/} -- logika biznesowa (walidacje, statusy, wykonanie),
    \item \texttt{data/} -- komunikacja z bazą danych,
    \item \texttt{model/} -- modele ORM i schematy odpowiedzi.
\end{itemize}

Najważniejsze modele w module to rekord powiadomienia i statusy/kanały. Warstwa \texttt{model} trzyma też schematy wejścia i wyjścia API, więc od razu wiadomo jakie pola przyjmujemy i co zwracamy.

W katalogu \texttt{service} są osobne pliki dla:
\begin{itemize}[leftmargin=1.5em]
    \item walidacji danych i czasu,
    \item maszyny stanów,
    \item dispatchera i delivery (faktyczna wysyłka),
    \item workera cyklicznego.
\end{itemize}

Taki podział jest wygodny do dalszego rozwoju, bo da się łatwo podmienić np. sposób wysyłki bez ruszania endpointów.

\section*{2. Interaktywna dokumentacja}

Dokumentacja jest zrobiona przez Swagger UI w FastAPI. Po uruchomieniu projektu można testować endpointy bezpośrednio z poziomu przeglądarki.

W dokumentacji widoczne są między innymi endpointy:
\begin{itemize}[leftmargin=1.5em]
    \item \texttt{POST /api/v1/notifications},
    \item \texttt{GET /api/v1/notifications},
    \item \texttt{GET /api/v1/notifications/\{id\}},
    \item \texttt{POST /api/v1/notifications/\{id\}/send-now},
    \item \texttt{POST /api/v1/notifications/process-ready},
    \item \texttt{GET /api/v1/metrics}.
\end{itemize}

Dzięki Swagger UI można od razu:
\begin{itemize}[leftmargin=1.5em]
    \item zobaczyć wymagane pola requestu,
    \item sprawdzić przykładowe odpowiedzi i kody HTTP,
    \item przetestować ręcznie scenariusze typu create -> send-now -> metrics.
\end{itemize}

To bardzo przyspiesza debugowanie, bo nie trzeba za każdym razem pisać ręcznie requestów.

\section*{3. Moduł planowania i obsługi powiadomień}

\subsection*{a) Architektura warstwowa (Controller -> Service -> Repository)}

Przepływ wygląda prosto:
\begin{enumerate}[leftmargin=1.8em]
    \item Controller odbiera żądanie HTTP.
    \item Service robi walidację i logikę biznesową.
    \item Repository zapisuje/odczytuje dane z bazy.
    \item Wynik wraca do klienta jako odpowiedź JSON.
\end{enumerate}

W praktyce wygląda to tak:
\begin{itemize}[leftmargin=1.5em]
    \item Controller nie zawiera logiki biznesowej, tylko mapuje błędy na kody HTTP (np. 404, 422).
    \item Service robi całą logikę: walidacja treści, walidacja odbiorcy, sprawdzanie czasu, ustawianie statusów.
    \item Repository ma proste funkcje typu \texttt{add}, \texttt{get\_by\_id}, \texttt{get\_ready}, \texttt{count\_by\_status}.
\end{itemize}

Taki układ ogranicza bałagan i ułatwia testowanie warstw osobno.

\subsection*{b) Maszyna stanów}

Powiadomienie może mieć statusy: \texttt{PENDING}, \texttt{SENT}, \texttt{FAILED}, \texttt{CANCELLED}.

W serwisie jest sprawdzanie dozwolonych przejść statusów. Dzięki temu nie da się przypadkowo zrobić niedozwolonej zmiany.

Przykładowe zasady:
\begin{itemize}[leftmargin=1.5em]
    \item wysyłkę (\texttt{send-now}) można wykonać tylko dla \texttt{PENDING},
    \item jeżeli delivery przejdzie poprawnie, status zmienia się na \texttt{SENT},
    \item jeżeli delivery rzuci wyjątek, status przechodzi na \texttt{FAILED},
    \item rekordów już wysłanych/anulowanych nie można ponownie wysłać.
\end{itemize}

To zabezpiecza system przed „skakaniem” po statusach w sposób nielogiczny.

\subsection*{c) Kanały komunikacji}

\subsubsection*{i. Powiadomienia PUSH}

Kanał PUSH działa przez wysłanie HTTP POST do lokalnego \texttt{ntfy}. Dispatcher wybiera kanał na podstawie pola \texttt{channel}.

\texttt{Ntfy.sh} to proste narzędzie do testowania powiadomień push. Dla laboratorium jest wygodne, bo od razu widać dostarczone wiadomości.

Implementacja w skrócie:
\begin{itemize}[leftmargin=1.5em]
    \item service pobiera rekord i wywołuje dispatcher,
    \item dispatcher dla \texttt{PUSH} wywołuje funkcję delivery,
    \item delivery robi request do \texttt{http://localhost:8088/\{topic\}}.
\end{itemize}

\begin{figure}[H]
    \centering
    % \includegraphics[width=0.9\textwidth]{ntfy_screenshot.png}
    \caption{Placeholder: zrzut z ntfy (zaplanowane i dostarczone powiadomienie + Swagger UI).}
\end{figure}

\subsubsection*{ii. Powiadomienia Email}

Kanał email działa przez MailHog (lokalny SMTP). Serwis delivery buduje wiadomość i wysyła ją na \texttt{localhost:1025}.

\texttt{MailHog} pozwala podejrzeć wysłane maile lokalnie, bez używania prawdziwej skrzynki.

Implementacja w skrócie:
\begin{itemize}[leftmargin=1.5em]
    \item dispatcher wykrywa kanał \texttt{EMAIL},
    \item delivery tworzy wiadomość (subject, from, to, content),
    \item wiadomość trafia do MailHog, gdzie można ją od razu podejrzeć w UI.
\end{itemize}

\begin{figure}[H]
    \centering
    % \includegraphics[width=0.9\textwidth]{mailhog_screenshot.png}
    \caption{Placeholder: zrzut z MailHog (dostarczone zaplanowane powiadomienie + Swagger UI).}
\end{figure}

\subsection*{d) Worker cykliczny}

Worker działa w tle równolegle do API (osobny wątek uruchamiany razem z aplikacją). Co kilka sekund pobiera rekordy gotowe do realizacji i je przetwarza. Dzięki temu powiadomienia mogą wysyłać się automatycznie.

Cykl działania workera:
\begin{enumerate}[leftmargin=1.8em]
    \item otwarcie sesji DB,
    \item pobranie rekordów \texttt{PENDING} z \texttt{scheduled\_at <= teraz\_utc},
    \item próba realizacji każdego rekordu,
    \item zapis nowego statusu,
    \item zamknięcie sesji i sleep (np. 5 sekund).
\end{enumerate}

W projekcie worker startuje razem z aplikacją, więc nie trzeba go uruchamiać ręcznie osobnym poleceniem.

\subsection*{e) Planowanie czasowe}

Czas zaplanowanej wysyłki jest konwertowany do UTC i porównywany z bieżącym czasem UTC (logika jak dla \texttt{DateTime.UtcNow}). Dzięki temu strefy czasowe użytkowników nie psują logiki harmonogramu.

Najważniejsze zasady:
\begin{itemize}[leftmargin=1.5em]
    \item przy wejściu przyjmujemy datę lokalną + nazwę strefy,
    \item przed zapisem konwertujemy do UTC i taką wartość trzymamy w DB,
    \item walidacja odrzuca datę z przeszłości,
    \item worker i metody wykonawcze porównują tylko względem UTC.
\end{itemize}

Dzięki temu niezależnie od lokalnej strefy użytkownika logika „czy rekord jest gotowy do wysyłki” działa spójnie.

\subsection*{f) Idempotencja}

Idempotencja jest potrzebna, żeby nie tworzyć podwójnych powiadomień przy ponowieniu żądania (np. po błędzie sieci).

W implementacji każde powiadomienie ma unikalny \texttt{idempotency\_key}. Klucz jest trzymany w bazie jako wartość unikalna.

Korzyści:
\begin{itemize}[leftmargin=1.5em]
    \item łatwiej wykryć duplikaty,
    \item można bezpieczniej ponawiać operacje po problemie sieciowym,
    \item łatwiej śledzić konkretną operację w logach i debugowaniu.
\end{itemize}

W projekcie klucz jest generowany po stronie backendu i zwracany w odpowiedzi API.

\section*{4. Prometheus}

Prometheus zbiera metryki z endpointu \texttt{/api/v1/metrics}. Dzięki temu można obserwować np. liczbę rekordów \texttt{PENDING}, \texttt{SENT}, \texttt{FAILED}, \texttt{CANCELLED} oraz podział po kanałach.

W praktyce metryki pomagają szybko ocenić stan systemu:
\begin{itemize}[leftmargin=1.5em]
    \item czy rośnie liczba \texttt{FAILED},
    \item czy kolejka \texttt{PENDING} nie stoi zbyt długo,
    \item jak rozkłada się ruch między \texttt{PUSH} i \texttt{EMAIL}.
\end{itemize}

Po uruchomieniu dockera i aplikacji Prometheus odpytuje endpoint cyklicznie i pokazuje wyniki w zakładkach \texttt{Table} i \texttt{Graph}.

\begin{figure}[H]
    \centering
    % \includegraphics[width=0.95\textwidth]{prometheus_chart.png}
    \caption{Placeholder: wykres metryk i statystyk w Prometheus.}
\end{figure}

\end{document}
